\chapterimage{figs/chapterhead.png}
\chapter{Applications, tools, exemplos em \texttt{C++}, testes e evolução do projeto}
\label{chap:applications_tools_exemplos_testes_evolucao}

\section{Introdução}
\label{sec:cap10_introducao}

Este capítulo encerra a estrutura principal deste manual do desenvolvedor ampliando novamente o foco. Nos capítulos anteriores, a atenção recaiu sobre a arquitetura central do sistema: organização do código-fonte, kernel, estrutura dos elementos do modelo, \textit{plugins} e parser. Agora, o objetivo é examinar como essa infraestrutura se projeta em aplicações concretas, ferramentas auxiliares, exemplos programáticos em \texttt{C++} e mecanismos de teste. Em outras palavras, este capítulo trata do ecossistema de desenvolvimento que se forma ao redor do núcleo do \textit{GenESyS}.

A importância desse fechamento é grande. Um simulador extensível não vive apenas de suas classes centrais. Ele também depende de aplicações que o exponham ao usuário, de ferramentas que apoiem seu uso e sua manutenção, de exemplos que mostrem como a biblioteca pode ser utilizada programaticamente e de testes que sustentem a evolução segura do sistema. Sem esses elementos, o núcleo do software até poderia existir, mas o projeto como plataforma de desenvolvimento ficaria empobrecido.

Por isso, este capítulo deve ser lido como uma síntese final da segunda metade do livro. Ele mostra como o GenESyS transita:
\begin{itemize}
    \item do kernel para as aplicações;
    \item da arquitetura para o uso programático;
    \item da extensão funcional para a manutenção técnica;
    \item da infraestrutura estável para a evolução futura do projeto.
\end{itemize}

\section{Applications como projeções do kernel}
\label{sec:cap10_applications_como_projecoes_kernel}

A pasta \texttt{source/applications} é o lugar onde o GenESyS deixa de aparecer apenas como infraestrutura e passa a assumir formas concretas de interação. Essa observação é fundamental para o desenvolvedor. Ela mostra que o projeto não foi desenhado para existir apenas como biblioteca abstrata, mas para sustentar aplicações reais sobre o seu núcleo.

Arquiteturalmente, isso é muito valioso. Ao separar o kernel das aplicações, o sistema ganha:
\begin{itemize}
    \item modularidade;
    \item reuso do núcleo em diferentes interfaces;
    \item menor acoplamento entre lógica de simulação e forma de interação;
    \item maior capacidade de evolução.
\end{itemize}

Na prática, essa separação permite que GUI, aplicação terminal e outras camadas futuras existam como projeções distintas do mesmo conjunto de serviços centrais.

\subsection{A GUI como aplicação sobre o núcleo}
\label{subsec:cap10_gui_como_aplicacao_sobre_nucleo}

Na primeira metade deste livro, a GUI foi tratada como o caminho principal de uso do simulador. Agora, ela passa a ser vista de outro modo: como aplicação construída sobre a infraestrutura do GenESyS. Essa mudança de perspectiva é especialmente importante para o desenvolvedor.

A GUI não é o núcleo. Ela é uma aplicação que:
\begin{itemize}
    \item consome o modelo;
    \item consome a semântica dos componentes e dados;
    \item se conecta aos mecanismos de simulação;
    \item expõe essas capacidades ao usuário por meio de uma interface visual.
\end{itemize}

Essa separação é arquiteturalmente correta e explica por que o sistema pode, em princípio, sustentar outras aplicações além da interface gráfica.

\subsection{A aplicação terminal como forma alternativa de uso}
\label{subsec:cap10_aplicacao_terminal_forma_alternativa}

A árvore de build atual do projeto mostra que a aplicação terminal é tratada de forma particularmente estruturada. O \texttt{CMakeLists.txt} de \texttt{source/applications/terminal} explicita o uso de uma aplicação chamada \texttt{GenesysShell}, bem como a possibilidade de selecionar outras classes de aplicação de terminal e até arquivos de exemplo específicos em \texttt{source/applications/terminal/examples} para compor o executável terminal :contentReference[oaicite:0]{index=0}

Isso é extremamente importante para o desenvolvedor. Mostra que o terminal não é apenas ferramenta auxiliar antiga ou caminho marginal de uso. Ele é uma aplicação efetiva do GenESyS e, mais do que isso, um ambiente valioso para:
\begin{itemize}
    \item estudar o uso programático do simulador;
    \item construir exemplos compactos;
    \item desenvolver aplicações leves sobre a biblioteca;
    \item verificar comportamentos sem a mediação da GUI.
\end{itemize}

\section{Os exemplos em \texttt{C++} como material de desenvolvedor}
\label{sec:cap10_exemplos_cpp_material_desenvolvedor}

Na primeira metade do livro, os exemplos em \texttt{C++} da aplicação terminal foram deliberadamente deixados de lado, porque o manual do usuário não deveria começar por código. Agora, porém, eles passam a ocupar um lugar central. Isso ocorre porque esses exemplos mostram uma das capacidades mais importantes do GenESyS como software: seu uso como biblioteca de simulação.

O \texttt{CMakeLists.txt} da aplicação terminal mostra explicitamente a existência de exemplos selecionáveis por build, inclusive exemplos \textit{smart} e exemplos mais elaborados de ensino, o que confirma o papel desses arquivos como material de estudo e demonstração arquitetural :contentReference[oaicite:1]{index=1}

\subsection{Por que esses exemplos são tão valiosos}
\label{subsec:cap10_por_que_exemplos_tao_valiosos}

Os exemplos em \texttt{C++} são valiosos porque permitem observar o simulador por outro ângulo. Na GUI, o modelo aparece como artefato visual e operacional. Nos exemplos programáticos, ele aparece como objeto construído explicitamente pelo desenvolvedor. Isso ajuda a revelar:
\begin{itemize}
    \item como o \textit{Simulator} é instanciado;
    \item como o modelo é construído por código;
    \item como componentes e \textit{data definitions} são criados e relacionados;
    \item como a execução é disparada;
    \item como a biblioteca do simulador pode ser reutilizada em novas aplicações.
\end{itemize}

\subsection{A ponte entre exemplos em \texttt{C++} e modelos salvos}
\label{subsec:cap10_ponte_exemplos_cpp_modelos_salvos}

Uma das características mais didaticamente ricas do projeto é a possibilidade de relacionar:
\begin{itemize}
    \item exemplos construídos programaticamente em \texttt{C++};
    \item modelos salvos na pasta \texttt{models/}.
\end{itemize}

Essa ponte é extremamente útil para o desenvolvedor porque mostra o mesmo sistema sob duas perspectivas:
\begin{itemize}
    \item como objeto de biblioteca;
    \item como artefato de modelagem e uso.
\end{itemize}

Essa correspondência ajuda a consolidar a compreensão do GenESyS como plataforma, e não apenas como aplicação pronta.

\section{Tools como camada de apoio}
\label{sec:cap10_tools_como_camada_apoio}

A pasta \texttt{source/tools} representa uma região importante do projeto, ainda que frequentemente menos comentada do que kernel, parser e \textit{plugins}. Sua função geral é apoiar uso, desenvolvimento, integração ou manutenção do sistema.

Do ponto de vista do desenvolvedor, ferramentas são relevantes por três razões:
\begin{itemize}
    \item ajudam a tornar o ecossistema do simulador mais produtivo;
    \item oferecem pontos de suporte à observação ou automação;
    \item podem servir como base para novas integrações.
\end{itemize}

Em termos arquiteturais, a existência de \texttt{tools} mostra que o projeto não está limitado à dicotomia ``kernel versus aplicação''. Há também espaço para instrumentos intermediários, utilitários e mecanismos de apoio.

\subsection{Quando estudar \texttt{tools}}
\label{subsec:cap10_quando_estudar_tools}

Nem sempre as ferramentas devem ser a primeira prioridade de leitura. Em geral, faz mais sentido estudá-las depois que o desenvolvedor já compreendeu:
\begin{itemize}
    \item o núcleo do simulador;
    \item o sistema de \textit{plugins};
    \item o parser;
    \item ao menos uma aplicação concreta, como a GUI ou a terminal.
\end{itemize}

Depois disso, a leitura das ferramentas tende a ser muito mais produtiva, porque o desenvolvedor consegue enxergar melhor sua função no ecossistema do projeto.

\section{A importância dos testes}
\label{sec:cap10_importancia_dos_testes}

Em um projeto do porte e da ambição do \textit{GenESyS}, a existência de testes é parte essencial da arquitetura evolutiva do software. Isso não significa apenas que o projeto ``tem testes'', mas que o próprio sistema de build já os trata como parte estruturada do desenvolvimento.

O \texttt{CMakeLists.txt} de \texttt{source/tests} mostra explicitamente a divisão entre testes unitários e \textit{smoke tests}, com um alvo agregado \texttt{genesys\_tests} e a possibilidade de habilitar ou desabilitar a construção dos testes de fumaça por opção de build :contentReference[oaicite:2]{index=2}

Essa organização é particularmente valiosa porque dá ao desenvolvedor um modo de pensar a validação em camadas:
\begin{itemize}
    \item testes unitários para comportamento local e invariantes menores;
    \item \textit{smoke tests} para sanidade mais ampla do sistema.
\end{itemize}

\subsection{Por que testes são parte da arquitetura e não só da rotina}
\label{subsec:cap10_testes_parte_arquitetura}

É comum pensar testes apenas como atividade externa ao projeto principal. No GenESyS, porém, eles devem ser lidos de forma mais forte: como parte da infraestrutura de evolução do sistema. Isso é particularmente importante porque o software tem:
\begin{itemize}
    \item kernel relativamente denso;
    \item mecanismos de extensão;
    \item múltiplas aplicações;
    \item integração entre parser, modelo, eventos e observabilidade.
\end{itemize}

Em um contexto assim, mudar código sem apoio de testes tende a ser arriscado. Portanto, o desenvolvedor deve ver a pasta \texttt{source/tests} como aliada direta na manutenção do projeto.

\subsection{Como o desenvolvedor deve usar os testes}
\label{subsec:cap10_como_usar_testes}

Uma prática madura de desenvolvimento no GenESyS deve considerar os testes em pelo menos três momentos:
\begin{enumerate}
    \item antes da alteração, para compreender a base atual de comportamento;
    \item durante a alteração, para evitar regressões locais;
    \item depois da alteração, para validar integração e sanidade geral.
\end{enumerate}

Essa disciplina é especialmente importante quando a mudança atinge kernel, parser, \textit{plugins} ou aplicações.

\section{Build, aplicações e seleção de comportamentos}
\label{sec:cap10_build_aplicacoes_selecao_comportamentos}

Ainda que o build não seja o eixo norteador do livro, ele continua oferecendo ao desenvolvedor um conjunto importante de sinais sobre a organização do projeto. O arquivo de build da aplicação terminal, por exemplo, mostra que o sistema foi pensado para permitir seleção configurável do comportamento do executável terminal, seja por classes específicas, seja por exemplos concretos em \texttt{source/applications/terminal/examples} :contentReference[oaicite:3]{index=3}

Isso revela algo importante sobre o projeto:
\begin{quote}
    O GenESyS não é apenas um programa; ele é uma plataforma sobre a qual comportamentos de aplicação podem ser escolhidos e montados.
\end{quote}

Essa observação interessa muito ao desenvolvedor porque reforça a modularidade do sistema e ajuda a pensar futuras aplicações sobre o núcleo.

\section{Evolução do projeto e caminhos de contribuição}
\label{sec:cap10_evolucao_projeto_caminhos_contribuicao}

A leitura combinada de kernel, parser, \textit{plugins}, applications, tools e tests mostra que o GenESyS é um projeto em evolução contínua. Essa evolução pode ocorrer em várias frentes:
\begin{itemize}
    \item aperfeiçoamento do kernel;
    \item ampliação do conjunto de componentes e \textit{data definitions};
    \item evolução da linguagem e do parser;
    \item fortalecimento da GUI e de outras aplicações;
    \item ampliação de ferramentas auxiliares;
    \item consolidação da base de testes.
\end{itemize}

Do ponto de vista do desenvolvedor, a pergunta mais útil deixa de ser:
\begin{quote}
    Onde eu posso mexer?
\end{quote}
e passa a ser:
\begin{quote}
    Em que camada do sistema minha alteração faz mais sentido arquiteturalmente?
\end{quote}

Essa mudança de perspectiva é o que distingue uma modificação oportunista de uma contribuição tecnicamente madura.

\subsection{Como escolher um ponto de entrada para contribuir}
\label{subsec:cap10_como_escolher_ponto_entrada}

A escolha do ponto de entrada depende muito do tipo de contribuição desejada.

\begin{itemize}
    \item Se a intenção for alterar a dinâmica central da simulação, o estudo deve começar no kernel.
    \item Se a intenção for ampliar linguagem e expressões, o estudo deve começar no parser.
    \item Se a intenção for acrescentar elementos de modelagem, a leitura deve concentrar-se em \textit{plugins}, componentes e \textit{model data}.
    \item Se a intenção for melhorar a experiência de uso, a GUI e as aplicações tornam-se centrais.
    \item Se a intenção for aumentar robustez, testes e tracing passam a ser prioridade.
\end{itemize}

Essa forma de pensar melhora bastante a qualidade da contribuição ao projeto.

\subsection{O papel dos exemplos e testes na evolução}
\label{subsec:cap10_papel_exemplos_testes_evolucao}

Exemplos e testes desempenham papéis complementares na evolução do software.

\begin{itemize}
    \item Os exemplos mostram \emph{como usar} o sistema e ajudam a explorar novas capacidades.
    \item Os testes ajudam a verificar \emph{se o sistema continua correto} ao ser modificado.
\end{itemize}

Um desenvolvedor maduro do GenESyS precisa valorizar ambos. Sem exemplos, o sistema perde capacidade de demonstração e aprendizado. Sem testes, perde segurança evolutiva.

\section{Uma visão final integrada do GenESyS}
\label{sec:cap10_visao_final_integrada_genesys}

Chegando ao final desta versão do manual do desenvolvedor, é útil voltar à visão integrada do projeto.

O \textit{GenESyS} pode ser compreendido como:

\begin{itemize}
    \item um kernel de simulação;
    \item uma hierarquia de objetos que representam modelo, componentes e dados;
    \item um sistema de linguagem e parser;
    \item uma infraestrutura de extensibilidade por \textit{plugins};
    \item um conjunto de aplicações construídas sobre o núcleo;
    \item um ecossistema de exemplos, ferramentas e testes que sustenta uso e evolução.
\end{itemize}

Essa visão integrada é importante porque evita leituras parciais demais. O GenESyS não é apenas GUI, não é apenas parser, não é apenas \textit{plugin framework}, nem apenas biblioteca de simulação. Ele é a articulação de todas essas dimensões.

\begin{table}[htbp]
    \centering
    \caption{Dimensões principais do GenESyS como plataforma de software.}
    \label{tab:cap10_dimensoes_principais_genesys}
    \small
    \begin{tabular}{|p{3.8cm}|p{7.4cm}|}
        \hline
        \textbf{Dimensão} & \textbf{Papel no projeto} \\
        \hline
        Kernel & Infraestrutura central da simulação \\
        \hline
        Parser & Linguagem, expressões e semântica textual do modelo \\
        \hline
        Plugins & Extensibilidade do simulador e ampliação dos elementos de modelagem \\
        \hline
        Applications & Formas concretas de uso, como GUI e aplicação terminal \\
        \hline
        Tools & Apoio ao uso, integração e manutenção \\
        \hline
        Tests & Segurança evolutiva, validação e prevenção de regressões \\
        \hline
        Examples & Material de demonstração, estudo e uso programático da biblioteca \\
        \hline
    \end{tabular}
\end{table}

\section{Considerações Finais}
\label{sec:cap10_consideracoes_finais}

Este capítulo fechou a estrutura principal do manual do desenvolvedor mostrando que o \textit{GenESyS} não deve ser lido apenas por suas classes centrais, mas também por seu ecossistema de aplicações, ferramentas, exemplos e testes. A GUI aparece agora como aplicação sobre o núcleo. A aplicação terminal e seus exemplos em \texttt{C++} revelam o uso do simulador como biblioteca. As ferramentas ampliam o ambiente de desenvolvimento. A árvore de testes fornece base para evolução segura. E a organização geral do projeto mostra que o GenESyS deve ser pensado como plataforma de software, e não apenas como simulador pronto.

Com isso, o leitor já dispõe de uma visão suficientemente ampla e tecnicamente fundamentada para usar, estudar e ampliar o GenESyS de forma madura. A partir daqui, o aprofundamento pode seguir diferentes trilhas: revisão e melhoria da GUI, estudo mais detalhado de componentes e \textit{plugins}, evolução do parser, fortalecimento da infraestrutura de testes, desenvolvimento de novas aplicações ou criação de novos domínios de modelagem sobre o mesmo núcleo. Em todos esses casos, a lógica permanece a mesma: compreender a camada correta, respeitar a arquitetura do sistema e fazer a evolução do simulador a partir de seus próprios princípios estruturais.

Teste seu entendimento desse conteúdo respondendo as seguintes perguntas:

\begin{enumerate}
    \item Por que a GUI deve ser entendida, do ponto de vista do desenvolvedor, como aplicação sobre o núcleo e não como o sistema inteiro?

    \item Qual é o valor arquitetural dos exemplos em \texttt{C++} da aplicação terminal?

    \item Por que a base de testes do projeto deve ser tratada como parte da infraestrutura de evolução do software?

    \item Como a distinção entre kernel, parser, \textit{plugins}, applications, tools, examples e tests ajuda o desenvolvedor a escolher melhor o ponto de entrada de uma contribuição?
\end{enumerate}